\subsection{计算结果与分析}

题目未给出特定的检测点数据，故构造合成算例检验算法。预设干扰源位置为 $G=(150,-250)$，根据检测点与 $G$ 的位置关系生成精确示向度，定位模型仍取 $\pm1^\circ$ 的误差范围。$G$ 仅用于生成验证数据，求解时只输入检测点坐标与示向度；具体输入见支撑材料中的算例文件。

\begin{table}[htbp]
    \centering
    \caption{典型布设下的定位区域直径}
    \label{tab:p1-results-concise}
    \begin{tabular}{lrrr}
        \toprule
        检测布设 & 检测点数 $k$ & 顶点数 $h$ & 直径 $D$/m \\
        \midrule
        双站正交交会 & 2 & 4 & 59.262511 \\
        在正交双站基础上增加第三站 & 3 & 6 & 59.262511 \\
        双站小交会角 $5^\circ$ & 2 & 4 & 1190.734431 \\
        双站同方位 & 2 & --- & 无有限直径 \\
        \bottomrule
    \end{tabular}
\end{table}

表\ref{tab:p1-results-concise}中各有界算例的旋转卡壳结果均与全部顶点对穷举结果一致。在1000个随机凸多边形的对照测试中，两种方法的最大数值差同样为0，表明当前测试范围内的直径计算结果一致；该结果用于核验实现，不替代理论论证，也不表示浮点计算相对于精确几何结果没有误差。

不同布设的交会结果存在明显差异。新增第三站后，区域顶点数由4增至6，但原最远点对仍被保留，因此直径未变。小交会角算例的直径显著增大，同方位算例则出现无界区域，说明定位不确定性与检测点的空间布设有关。由于小交会角算例同时改变了一个检测距离，上述比较不单独用于量化交会角的影响。

进一步固定已有测向数据，逐个增加检测点。此时
\begin{equation}
    P_{k+1}=P_k\cap W_{k+1}\subseteq P_k,
    \qquad D(P_{k+1})\le D(P_k),
\end{equation}
其中假定交集保持非空有界。在同一组检测点的嵌套实验中，检测点数由2增至12时，直径由69.309097 m降至46.408654 m，降低约33.04\%；10次新增观测中有4次直径在报告精度下保持不变。因此，增加检测点可以保证直径不增，却不保证严格减小；新观测能否进一步缩小直径，取决于其是否排除了原区域的最远点对，而非仅取决于检测点数量。
